\chapter{计算系统全景}

第一册已经说明程序如何变成进程、虚拟地址、调用约定和机器指令。第二册继续追问：当一个程序需要利用多个核心、多个内存层级、多个线程甚至多台机器时，计算能力从哪里来，又会在哪里损失。这个问题不能只从代码出发，也不能只从硬件手册出发。真正的计算系统由硬件、操作系统、运行时、算法和通信协议共同构成。

\section{从一条指令到一个系统}

一条加法指令看起来很简单，但它并不是孤立执行的。核心要取指、译码、重命名寄存器、等待操作数、选择执行端口、写回结果并提交状态。如果操作数在寄存器里，成本很低；如果来自 L1 cache，延迟增加；如果来自 L3、远端 NUMA 节点或磁盘，程序看到的时间会跨越几个数量级。并发程序还要考虑另一个核心是否正在写同一条 cache line，分布式程序还要考虑消息是否跨越网络。

因此本册把计算系统分成五层。第一层是硬件执行，包括核心、缓存、TLB、内存控制器和互连。第二层是操作系统，包括调度、页表、系统调用、I/O 和计数器。第三层是运行时，包括线程池、任务队列、同步原语和内存分配器。第四层是算法，包括归约、扫描、分区、图遍历和负载均衡。第五层是分布式协议，包括 RPC、复制、共识、shuffle 和检查点。

\begin{keyidea}
高性能不是某一层单独决定的。硬件给出能力边界，操作系统分配资源，运行时组织执行，算法决定数据依赖，分布式协议处理通信和故障。
\end{keyidea}

\section{计算成本的基本单位}

分析性能时，不能只看“执行了多少行代码”。更稳定的单位是指令数、周期数、访存字节数、cache miss、TLB miss、分支误预测、同步等待时间、上下文切换、系统调用次数和网络往返延迟。不同单位之间不能随便相加，但它们共同解释程序为什么慢。

例如，一个数组求和循环的算术量很小，主要成本是把数据从内存送到核心。一个矩阵乘如果分块得当，会重复使用缓存中的数据，瓶颈可能转向浮点执行端口。一个共享队列如果每次 push 都抢同一把锁，瓶颈不是算术，也不是内存带宽，而是临界区串行化和 cache line 所有权迁移。

\section{延迟、吞吐和并发度}

系统分析经常混淆三个词：延迟、吞吐和并发度。延迟是单个操作从开始到结束的时间，吞吐是单位时间完成多少操作，并发度是同时在系统中飞行的操作数量。一个内存 load 的延迟可能很高，但如果核心能同时挂起多个 miss，整体吞吐仍然可以接受；一个网络请求延迟很高，但如果客户端保持足够多的 in-flight 请求，吞吐可以接近带宽上限。

这三个量之间有一个朴素但非常有用的关系：要填满一个延迟很高的管道，需要足够并发度。磁盘、网络、内存和 CPU 后端都可以用这个心智模型理解。区别在于，CPU 的并发度来自乱序窗口、load buffer、SIMD lane 和多核心；网络系统的并发度来自连接数、请求队列和异步 I/O；分布式计算的并发度来自分片、任务和流水线。

\begin{mentalmodel}
延迟回答“一个操作要等多久”，吞吐回答“系统每秒能做多少”，并发度回答“为了达到吞吐，需要同时容纳多少未完成工作”。
\end{mentalmodel}

\section{Amdahl 和 Gustafson}

多核性能讨论离不开可并行比例。Amdahl 定律说明，如果程序有一部分必须串行执行，那么线程数增加到一定程度后，加速会被串行部分限制。Gustafson 定律从另一个角度看问题：如果随着核心数增加而扩大问题规模，可并行部分也随之增加，系统仍然能获得更高总吞吐。

这两个定律不是考试公式，而是工程判断工具。一个服务端请求路径中，如果日志锁、全局分配器、单分片元数据或单线程事件循环占比很高，增加 worker 不会线性加速。一个离线计算任务如果可以把输入切成更多分片，让每个节点处理独立数据，再做树形合并，扩大规模时就更接近 Gustafson 的场景。

\section{观察系统的三张图}

学习计算系统时，建议每遇到一个程序都画三张图。第一张是数据流图：输入从哪里来，经过哪些结构，输出到哪里。第二张是执行图：哪些线程、任务或节点执行哪些阶段，哪里并行，哪里等待。第三张是成本图：每个阶段主要消耗 CPU、内存、I/O、锁等待还是网络。

这三张图可以避免把问题混成一团。一个队列积压，可能是消费者计算慢，也可能是生产者太快，也可能是下游 I/O 阻塞，也可能是队列锁竞争。只有先画清数据和执行边界，测量才知道该放在哪里。

\section{正确性先于性能}

并发和分布式计算最容易出现的错误，是把速度放在正确性之前。数据竞争可能让程序偶尔错；错误的内存序可能只在某些 CPU 上触发；无锁结构如果没有处理对象回收，可能在压力测试里才出现悬垂访问；分布式系统如果没有处理超时和重试，可能在网络抖动时重复提交请求。

本册每次讨论优化，都要先给出语义边界。共享变量由谁写，谁读，什么时候发布，什么时候回收。任务是否可以重排，结果是否依赖顺序。消息是否幂等，故障后是否可以重试。只有这些问题说清楚，性能数字才有意义。

\section{本册贯穿实验}

第二册的实验项目叫 Compute Systems Labs。第一阶段实现顺序归约、并行归约、带锁计数器、原子计数器和有界队列。后续章节会围绕它继续扩展：加入线程池、任务图、工作窃取、无锁队列、并行扫描、I/O 背压和分布式模拟。

\begin{labbox}
实验一：运行 Compute Systems Labs。记录 CPU 型号、核心数、线程数、缓存层级、系统版本、编译器版本和测试命令。不要只写测试通过，还要解释每个模块分别验证了计算系统的哪一类问题。
\end{labbox}
